% =============================================================================
% conclusion.tex — Заключение
% =============================================================================

\section*{ЗАКЛЮЧЕНИЕ}
\addcontentsline{toc}{section}{ЗАКЛЮЧЕНИЕ}

В рамках настоящей работы была разработана и реализована система автоматической генерации технической документации на основе графовой модели кода и больших языковых моделей. Основные результаты работы:

\begin{enumerate}
  \item Проведён анализ существующих подходов к автоматической генерации документации (code summarization, LLM в IDE, Swagger/OpenAPI, RAG, Code Graph + LLM) и выявлены их ключевые ограничения: локальность контекста, невозможность анализа межсервисного взаимодействия, зависимость от облачных провайдеров.

  \item Разработана унифицированная графовая модель представления кода, включающая 24~типа узлов и 24~типа связей, покрывающая структурные, вызовные и интеграционные отношения. Модель обеспечивает представление кодовой базы на уровне, достаточном для генерации документации различной степени детализации.

  \item Спроектирована модульная архитектура системы на основе принципов DDD с выделением 8~ограниченных контекстов (git, graph, chunking, embedding, ai, rag, library, postprocess) и 3~модулей ядра (domain, db, shared), обеспечивающая расширяемость и тестируемость.

  \item Реализован конвейер обработки данных из 7~этапов: клонирование репозитория, построение графа кода (PSI-парсинг), линковка связей, анализ библиотек (ASM), семантическое чанкование, векторизация и многоступенчатая генерация документации.

  \item Разработана подсистема RAG с графовым расширением контекста, комбинирующая точный поиск по FQN, обход графа зависимостей и векторный поиск для формирования контекста ответа.

  \item Реализована визуализация графа кода в двух режимах: анализ одного приложения и кросс-приложенческое представление интеграций (HTTP, Kafka, Camel).

  \item \todo{Привести результаты тестирования и оценки качества: 107 тестовых файлов, покрытие всех 8 контекстов, интеграционные тесты с TestContainers (PostgreSQL 16 + pgvector), результаты апробации на реальных проектах (количество узлов/рёбер, время индексации, Cyrillic ratio docPublic, precision/recall RAG-ответов, сравнение наивного RAG vs графового расширения).}
\end{enumerate}

\textbf{Направления дальнейшего развития:}
\begin{itemize}
  \item расширение поддержки языков программирования (Java-парсинг на уровне PSI);
  \item добавление кеширования для больших наборов приложений;
  \item реализация инкрементального обновления графа при коммитах;
  \item разработка метрик архитектурной полноты документации;
  \item экспорт документации в различные форматы (Confluence, PDF, Markdown).
\end{itemize}
